\begin{solapartat}
\begin{minipage}[t]{0.55\linewidth}
Veient que $R \gg h$ i que, tal com es mostra a la gràfica, es considera que $E_\mathrm{p} = 0$ a la superfície del planeta. Podem expressar l'energia potencial gravitatòria com: $[E_p = mgh]$ en què $g$ és l'acceleració de la gravetat a la superfície del planeta.
\end{minipage}\hfill
\begin{minipage}[t]{0.42\linewidth}
\vspace{0pt}
\figura[1]{sol_fig1.png}
\end{minipage}

\medskip
\punts{0,2} $\mathit{pendent} = \dfrac{\Delta y}{\Delta x} = \dfrac{\Delta E_p}{\Delta h} = \dfrac{80}{20} = 4\ \mathrm{J/m}$

\medskip
\punts{0,4}
$\begin{aligned}[t]
&E_p = mgh\\
&\mathit{pendent} = mg \Rightarrow g = \frac{\mathit{pendent}}{m} = \frac{4}{2} = 2\ \mathrm{m/s^2}
\end{aligned}$

\medskip
\punts{0,2} $g = \dfrac{GM}{R^2} \Rightarrow M = \dfrac{gR^2}{G}$

\medskip
\punts{0,2} $M = \dfrac{2\cdot\left(5\times 10^{6}\right)^2}{6,67\times 10^{-11}} = 7,50\times 10^{23}\ \mathrm{kg}$
\end{solapartat}

\begin{solapartat}
\punts{0,2} La deducció es basa en el principi de conservació de l'energia mecànica. L'energia mecànica a l'infinit és zero.

\medskip
\punts{0,6}
$\begin{aligned}[t]
&E_{\text{mecànica inicial}} = E_{\text{mecànica final}}\\
&K_{\textit{inicial}} + U_{\textit{inicial}} = K_{\textit{final}} + U_{\textit{final}}\\
&\frac{1}{2}mv_{esc}^2 - G\frac{Mm}{R} = 0 + 0 \Rightarrow v_{esc} = \sqrt{\frac{2GM}{R}}
\end{aligned}$

\medskip
\punts{0,2} $v_{esc} = \sqrt{\dfrac{2\cdot 6,67\times 10^{-11}\cdot 7,50\times 10^{23}}{5\times 10^{6}}} = 4,47\times 10^{3}\ \mathrm{m/s}$
\end{solapartat}
